% DRAFT (Claude, 2026-11-27, Anders): proposal for the whole section "Deterministic and Stochastic
% Optimisation" after the supervisor's feedback of 27 November. The group's sentences are kept; lines
% marked NEW are added, lines marked CHANGED are reworded. To be reworked by the group.
% Pages in birgeLouveaux2011 (2nd ed.) behind each part are kept as "% src:" comments, since the
% in-text citations are removed as the supervisor asked.

\subsection{Deterministic and Stochastic Optimisation}
\label{subsec:theory_det_stoch}

% NEW (supervisor: cite the textbooks once)
The theory presented in this section is based on \textcite{lundgren2010} for linear programming and duality, \textcite{boydVandenberghe2004} for convexity, and \textcite{birgeLouveaux2011} for stochastic optimisation.

\subsubsection{Deterministic Optimisation}
\label{subsec:deterministic_optimisation}

Deterministic optimisation is a mathematical approach in which all model parameters are represented by fixed, known values. The objective is to find the best possible solution by maximising or minimising an objective function, subject to a set of constraints. The model determines the optimal values of the decision variables for the specified parameters. Since uncertainty is not explicitly represented, the resulting solution is optimal for the assumed input values.

\paragraph{Linear programming}
% CHANGED: "Lundgren et al. (2010) write an LP problem in ..." -> no in-text citation
If the objective function and all constraints are linear, and all variables are continuous, the problem is a linear programming (LP) problem. An LP problem can be written in the following general form:
\begin{equation}
\begin{array}{@{}l@{\hspace{3.2em}}l@{\;\;}l@{\quad}l@{}}
\min & \llap{$z={}$}{\displaystyle\sum_{j=1}^{n}}\, c_j x_j & & \\[4ex]
\text{s.t.} & {\displaystyle\sum_{j=1}^{n}}\, a_{ij} x_j & \le b_i, & i = 1,\ldots,m\\[3ex]
 & \multicolumn{1}{r@{\;\;}}{x_j} & \ge 0, & j = 1,\ldots,n
\end{array}
\label{eq:lp_general}
\end{equation}
Here, $x_j$ are the decision variables and $z$ is the value of the objective function. The parameter $c_j$ is the objective function coefficient for variable $x_j$, $a_{ij}$ is the constraint coefficient for variable $x_j$ in constraint $i$, and $b_i$ is the right-hand side of constraint $i$. The last line states that the variables cannot be negative. A maximisation problem can be written in the same form, since maximising $z$ is the same as minimising $-z$.

% REMOVED (Anders 2026-11-27): "Since the feasible region is convex and the objective function is linear, every local optimum of an LP problem is also a global optimum." Now covered by the convexity paragraph, which defines "convex" first.

% NEW (Anders 2026-11-27: a paragraph on convexity, since SDDP relies on it). Every statement is from boydVandenberghe2004; pages in the src comments.
\paragraph{Convexity}
A set is convex when it contains every point on the straight line between any two of its own points. A function $f$ defined on a convex set is convex when its value at any point between $x$ and $y$ is no larger than the corresponding weighted average of $f(x)$ and $f(y)$. For all $x$ and $y$ in the set and all $\theta$ with $0\le\theta\le 1$,
\begin{equation}
f\bigl(\theta x+(1-\theta)y\bigr)\le\theta f(x)+(1-\theta)f(y).
\label{eq:convex_function}
\end{equation}
A function $f$ is concave if $-f$ is convex. % src: boydVandenberghe2004 p. 23 (convex set), p. 67 (convex and concave function, eq. 3.1)
An optimisation problem is convex if a convex objective function is minimised over a convex set, and every LP problem is a convex problem. % src: boydVandenberghe2004 p. 137; p. 2
Two properties of convex problems are used later in this chapter. First, every local optimum of a convex problem is also a global optimum. % src: boydVandenberghe2004 p. 138 (Sec. 4.2.2)
Second, the linear approximation of a convex function at one point, given by its value and slope there, never lies above the function, so it gives a bound that holds everywhere. % src: boydVandenberghe2004 p. 69 (eq. 3.2 and the text after it; stated there for differentiable functions)
% MOVE TO THE SDDP SECTION (Anders 2026-11-27): the inequality f(y) >= f(x) + grad f(x)^T (y - x) (boydVandenberghe2004 p. 69) and "the maximum of a set of linear functions is a convex, piecewise linear function" (p. 80) belong where the cuts are introduced.

\paragraph{Nonlinear and integer problems}
% CHANGED: shortened to what the thesis uses (supervisor: the NLP/MIP part can be cut down)
Nonlinear programming (NLP) allows the objective function or constraints to be nonlinear. Such problems are not necessarily convex, and in a nonconvex problem a local optimum is not necessarily a global optimum. % src: boydVandenberghe2004 pp. 9--10 (nonconvex problems and local optimisation)
Integer programming (IP) restricts decision variables to integer values, and problems containing both continuous and integer variables are referred to as mixed-integer programming (MIP). A set of integer points is not convex by the definition above, since the point halfway between two neighbouring integers is not an integer. Problems with integer variables are therefore not convex. % own argument from the definition, no source needed
% MOVED (2026-11-28) to the SDDP section: SDDP requires convex stage problems (helseth2016 p. 938).

\paragraph{Dual problem}
Every linear programming problem has an associated dual problem, constructed from the coefficients of the original, or primal, problem. Each constraint in the primal problem corresponds to a variable in the dual problem.
% NEW: the dual written out
For the LP problem in \eqref{eq:lp_general}, with one dual variable $v_i$ for each constraint $i$, the dual problem is
\begin{equation}
\begin{array}{@{}l@{\hspace{3.2em}}l@{\;\;}l@{\quad}l@{}}
\max & \llap{$w={}$}{\displaystyle\sum_{i=1}^{m}}\, b_i v_i & & \\[4ex]
\text{s.t.} & {\displaystyle\sum_{i=1}^{m}}\, a_{ij} v_i & \le c_j, & j = 1,\ldots,n\\[3ex]
 & \multicolumn{1}{r@{\;\;}}{v_i} & \le 0, & i = 1,\ldots,m
\end{array}
\label{eq:lp_dual}
\end{equation}
% CHANGED: "These dual variables can be interpreted as shadow prices ..." made specific
The dual variables can be interpreted as shadow prices. $v_i$ is the change in the optimal objective value $z$ when the right-hand side $b_i$ increases by one unit. If the primal or dual problem has a bounded optimal solution, the two problems must have the same objective value at optimum. This property is known as strong duality. % src: boydVandenberghe2004 pp. 224--227

% MOVED (2026-11-28) to the Benders section: the dual variables of the linking constraints give the slope of the cut. The water value as the dual of the reservoir balance is in the background chapter.

\subsubsection{Stochastic Optimisation}
\label{subsec:stochastic_optimisation}

In many decision problems, some parameters are not known when decisions must be made. Stochastic optimisation incorporates this uncertainty into the mathematical model by representing uncertain parameters as random variables with probability distributions. % src: birgeLouveaux2011 pp. 55--58

In practice, uncertainty is often represented by a finite set of scenarios, where each scenario represents a possible realisation of the uncertain parameters that is associated with a probability. Rather than optimising for a single forecast, stochastic optimisation considers multiple possible outcomes. A common objective is to maximise expected profit, calculated as the probability-weighted average profit across scenarios. This is referred to as risk-neutral optimisation, since variability in outcomes is not explicitly penalised. % src: birgeLouveaux2011 p. 56; pp. 7--8, 66--67

\subsubsection{Two-Stage Stochastic Optimisation}
\label{subsec:two_stage_stochastic_optimisation}

A two-stage stochastic optimisation problem with recourse consists of two stages of decision-making. First-stage decisions are made before the uncertain parameters are known. These are called \textit{here-and-now} decisions and must be the same across all scenarios. Once the uncertainty is realised, second-stage decisions can be made to adapt to the observed outcome. These are called recourse decisions and may differ between scenarios. % src: birgeLouveaux2011 pp. 57--59, 164
First-stage decisions must account for the costs or benefits associated with subsequent decisions. The recourse function captures these consequences by expressing the expected optimal second-stage objective value for a given first-stage decision.

% NEW (supervisor: the two-stage problem and the recourse function as expressions)
Let $x$ denote the first-stage decisions, $y$ the second-stage decisions and $\xi$ the uncertain parameters. The two-stage stochastic linear program is
\begin{equation}
\begin{array}{@{}l@{\quad}l@{}}
\min & c^{\top}x + \mathcal{Q}(x)\\[1ex]
\text{s.t.} & Ax = b,\quad x \ge 0,
\end{array}
\label{eq:two_stage}
\end{equation}
where $c$, $A$ and $b$ are known first-stage data and $\mathcal{Q}(x)=\mathbb{E}_{\xi}\left[Q(x,\xi)\right]$ is the recourse function. For a given first-stage decision $x$ and realisation $\xi$, the second-stage value is
\begin{equation}
Q(x,\xi)=\min_{y}\left\{\,q(\xi)^{\top}y \;\middle|\; Wy = h(\xi)-T(\xi)x,\; y\ge 0\right\},
\label{eq:second_stage}
\end{equation}
where $q(\xi)$, $h(\xi)$ and $T(\xi)$ are second-stage data that become known with $\xi$, and $W$ is the same in all scenarios. % src: birgeLouveaux2011 eqs. (4.1)-(4.4), pp. 59-60
% CHANGED: "The objective therefore combines ..." now refers to the expression
The objective in \eqref{eq:two_stage} therefore combines the first-stage value with the expected second-stage value. A major computational challenge is evaluating the recourse function for different first-stage decisions. % src: birgeLouveaux2011 pp. 103--105
% NEW: the property the cuts rely on
With a finite number of scenarios, $\mathcal{Q}(x)$ is a piecewise linear and convex function of $x$. In a maximisation problem it is concave. % src: birgeLouveaux2011 Theorem 2 and the sentence after it, Sec. 3.1, p. 107
This property makes it possible to approximate the recourse function by linear cuts.

% CHANGED: wait-and-see and EVPI with expressions; link to the deterministic problem
The \textit{wait-and-see} problem differs from the recourse problem by assuming that the uncertain parameters are known before any decisions are made. Each scenario can therefore be optimised separately, including its first-stage decisions. For each scenario, this is a deterministic problem of the type described in Section~\ref{subsec:deterministic_optimisation}. Let $z(x,\xi)=c^{\top}x+Q(x,\xi)$ denote the objective value of the first-stage decision $x$ in scenario $\xi$. The optimal values of the recourse problem (RP) and the wait-and-see problem (WS) are
\begin{equation}
\mathrm{RP}=\min_{x}\,\mathbb{E}_{\xi}\left[z(x,\xi)\right],
\qquad
\mathrm{WS}=\mathbb{E}_{\xi}\left[\min_{x}\, z(x,\xi)\right].
\label{eq:rp_ws}
\end{equation}
Since the wait-and-see solution assumes perfect information about future outcomes, it provides an optimistic bound on the optimal expected objective value of the recourse problem. The difference between these two values is called the expected value of perfect information (EVPI),
\begin{equation}
\mathrm{EVPI}=\mathrm{RP}-\mathrm{WS},
\label{eq:evpi}
\end{equation}
which measures the maximum expected benefit of knowing the uncertain parameters in advance. % src: birgeLouveaux2011 eqs. (1.1)-(1.4), pp. 163--165
% NEW: sign for maximisation, and the link to the case study
In a maximisation problem, the wait-and-see value is the larger of the two and $\mathrm{EVPI}=\mathrm{WS}-\mathrm{RP}$.

\subsubsection{Multistage Stochastic Optimisation}
\label{subsec:multistage_stochastic_optimisation}

A multistage stochastic optimisation problem extends the two-stage framework to several stages. Information is revealed gradually between stages, and decisions at each stage are based on the information available at that time and the probability distributions of future uncertain outcomes. % src: birgeLouveaux2011 pp. 149--151
Decisions in one stage may affect decisions in later stages.

% NEW (supervisor: introduce notation that the DP and SDDP parts can build on)
With decisions $x_t$ in stage $t=1,\ldots,N$, the problem can be written as a recursion. For $t=2,\ldots,N$, the value of stage $t$ for a given previous decision $x_{t-1}$ and realisation $\xi_t$ is
\begin{equation}
\begin{array}{@{}l@{\quad}l@{}}
Q_t(x_{t-1},\xi_t)=\;\min & c_t(\xi_t)^{\top}x_t+\mathcal{Q}_{t+1}(x_t)\\[1ex]
\phantom{Q_t(x_{t-1},\xi_t)=\;}\text{s.t.} & W_t x_t = h_t(\xi_t)-T_{t-1}(\xi_t)x_{t-1},\quad x_t\ge 0,
\end{array}
\label{eq:multistage_recursion}
\end{equation}
where $\mathcal{Q}_{t+1}(x_t)=\mathbb{E}_{\xi_{t+1}}\left[Q_{t+1}(x_t,\xi_{t+1})\right]$ is the expected value of all later stages and $\mathcal{Q}_{N+1}=0$. The first stage minimises $c_1^{\top}x_1+\mathcal{Q}_2(x_1)$ subject to the first-stage constraints. % src: birgeLouveaux2011 eqs. (4.1)-(4.4), pp. 150--151
The previous decision $x_{t-1}$ is the only link between two stages and acts as the state of the system. If the distribution of $\xi_{t+1}$ depends on earlier outcomes, this information must also be part of the state. % src: birgeLouveaux2011 p. 151
% MOVED (2026-11-28): the link from this recursion to dynamic programming is to be made in the dynamic programming section (supervisor point on "stage").

A scenario describes one possible development of the uncertain parameters over the full planning horizon. The scenarios can be represented in a scenario tree, as illustrated in Figure~\ref{fig:MSSP}. Each node represents the information available at a particular stage, while the branches describe possible subsequent outcomes. Except for the root node, every node is connected to exactly one node in the preceding stage. % src: birgeLouveaux2011 pp. 22, 152
The solution therefore consists of decisions that depend on the observed history, rather than one fixed sequence of decisions. % src: birgeLouveaux2011 pp. 21--22
As the planning horizon grows, the number of possible scenario paths may increase rapidly, leading to substantial computational requirements. % src: birgeLouveaux2011 p. 152

\begin{figure}[H]
    \centering
    \includegraphics[width=\textwidth]{Images/MSSP.png}
    \caption{Illustration of two-stage and multistage stochastic optimisation using scenario trees.}
    \label{fig:MSSP}
\end{figure}

A decision stage is defined by the information available when decisions are made, whereas a time period represents an interval in the planning horizon. These two concepts do not necessarily coincide, since several time periods may be included within the same decision stage. % src: birgeLouveaux2011 pp. 58, 65
